\documentclass[11pt]{amsart}
\usepackage[T1]{fontenc}
\usepackage[utf8]{inputenc}
\usepackage[french]{babel}
\usepackage{amssymb,stmaryrd}
\usepackage[margin=1.05in]{geometry}
\usepackage[hidelinks]{hyperref}
\newcommand{\stamptime}{\the\numexpr\time/60\relax:\ifnum\numexpr\time-60*(\time/60)\relax<10 0\fi\the\numexpr\time-60*(\time/60)\relax}
\AddToHook{shipout/foreground}{\put(0.5\paperwidth,-\paperheight+0.55cm){\makebox[0pt]{\tiny\ttfamily unit/doc/\jobname.tex \ --- \ compil\'e le \today, \stamptime}}}
\newtheorem{theorem}{Théorème}
\newtheorem{proposition}[theorem]{Proposition}
\newtheorem{lemma}[theorem]{Lemme}
\newtheorem{corollary}[theorem]{Corolaire}
\theoremstyle{definition}
\newtheorem{definition}[theorem]{Définition}
\theoremstyle{remark}
\newtheorem{remark}[theorem]{Remarque}
\DeclareMathOperator{\Ann}{Ann}
\DeclareMathOperator{\Ker}{Ker}
\DeclareMathOperator{\Ima}{Im}
\DeclareMathOperator{\Gr}{Gr}
\renewcommand{\Pr}{\operatorname{Pr}}
\newcommand{\A}{\mathbf{A}}
\newcommand{\B}{\mathbf{B}}
\newcommand{\ua}{\underline{a}}
\newcommand{\fa}{\mathfrak{a}}
\newcommand{\fb}{\mathfrak{b}}
\newcommand{\iv}[1]{\llbracket #1\rrbracket}
\title{Remarques sur la section XX-9 d'ACMC3 :\\
profondeur homologique $\Pr=\Gr$ et le théorème XX-9.10}
\date{1\textsuperscript{er} novembre 2026}
\begin{document}
\maketitle

\noindent Les numéros renvoient à la version de \texttt{ACMC3.pdf}
(chapitre XX, \emph{Résolutions libres finies}) consultée le 1\textsuperscript{er}
novembre 2026. On note $K^\bullet(\ua,E)$ le complexe de Koszul montant,
$K_\bullet(\ua,E)$ le complexe descendant, $H^j$ et $H_j$ leurs homologies,
$\fa=\langle\ua\rangle$, $n$ la longueur de $\ua$.

Résumé.
\begin{enumerate}
\item Deux corrections préliminaires (indices dans 9.3 ; profondeur infinie
  et modules de type fini).
\item Une démonstration réellement directe de 9.7 n'utilisant que McCoy,
  9.4 et 9.6 ; elle fournit en corolaires 1.4, 1.6, 1.13, 1.14.
\item Une lecture de 9.8 via $\Pr$ (avec l'hypothèse « $E$ de présentation
  finie »).
\item Une démonstration complète de 9.10.
\end{enumerate}

\section{Deux corrections préliminaires}

\subsection{Indices dans la définition 9.3}
Le complexe montant commence en $K^0(\ua,E)=E$ et
$H^0(\ua,E)=\Ann_E(\fa)$. La définition doit être
\[
\Pr_\A(\ua,E)\ge k\iff \forall j\in\iv{0..k-1},\ H^j_\A(\ua,E)=0 .
\]
Avec $\iv{1..k}$, pour $n=1$ et $a$ régulier pour $E$ avec $E/aE\ne0$, on
aurait $\Gr=1$ mais $H^1=E/aE\ne0$. L'« Exemple » suivant 9.3 (si $\ua$
est de longueur $k$, $\Pr\ge k+1$ donne $H^k=E/\fa E=0$) et le lemme 9.6
confirment l'indexation à partir de $0$. De même, dans 9.4 il faut lire
« pour chaque $j\ge0$ ».

\subsection{Profondeur infinie}
L'« Exemple » de 9.3 (« si $\fa E=E$, la profondeur homologique est
infinie ») et la proposition 1.7.3 demandent que $E$ soit de type fini.

\emph{Contre-exemple.} $\A=\mathbb Z$, $\ua=(2)$, $E=\mathbb Q/\mathbb Z$ :
$2E=E$ mais $H^0=\Ann_E(2)\simeq\mathbb Z/2\ne0$, donc $\Pr=\Gr=0$.

\emph{Cas $E$ de type fini.} $\fa E=E$ donne (lemme du déterminant)
$x\in\fa$ avec $(1+x)E=0$. Les $H^j$ sont tués par $x$ (lemme 9.2) et par
$1+x$ (lemme 9.11), donc sont nuls.

\section{Le théorème 9.7 : une démonstration directe}

La démonstration du texte (« lemme 9.6, théorème 9.4, et $\Gr$ vérifie les
mêmes propriétés ») appelle deux remarques.
\begin{enumerate}
\item Le cas de base est implicite : $\Gr\ge1$ signifie qu'un polynôme de
  Kronecker $f$ est régulier pour $E[X]$, et $\Pr\ge1$ que
  $\Ann_E(\fa)=0$ ; l'équivalence est le lemme de McCoy pour les modules.
\item Elle utilise côté $\Gr$ les théorèmes 1.13 et surtout 1.14, qui sont
  des résultats profonds de RLF (1.14 repose sur 1.15). La démonstration
  n'est donc « directe » qu'en apparence.
\end{enumerate}
En fait on peut démontrer 9.7 en n'utilisant que McCoy, 9.4 et 9.6, et
obtenir ensuite les résultats de la section 1 comme corolaires.

\begin{lemma}[McCoy pour les modules]\label{mccoy}
Soient $f\in\A[T]$ et $E$ un $\A$-module. Si $\Ann_E(c(f))=0$, alors $f$
est régulier pour $E[T]$. La réciproque est immédiate.
\end{lemma}
\begin{proof}
On applique le lemme de Dedekind-Mertens III-2.1 dans l'anneau
$\A\ltimes E$ (idéalisation de Nagata, $E$ idéal de carré nul) : si
$g\in E[T]$, $fg=0$ et $m\ge\deg g$, on obtient $c(f)^{m+1}c(g)=0$. Une
puissance d'un idéal $E$-régulier est $E$-régulière (si $IJx=0$, alors
pour $j\in J$, $I(jx)=0$, donc $jx=0$, donc $Jx=0$, donc $x=0$). Donc
$c(g)=0$, i.e.\ $g=0$. Réciproquement si $\fa x=0$ avec $x\in E$, alors
$fx=0$.
\end{proof}

\begin{theorem}[9.7]
Pour tout $k\ge1$, $\Pr_\A(\ua,E)\ge k\iff\Gr_\A(\ua,E)\ge k$. Plus
précisément, chacune de ces conditions équivaut à : \emph{toute} suite de
Kronecker de longueur $k$ attachée à $\fa$ est $E$-régulière.
\end{theorem}

\begin{proof}
Soit $(f_1,\dots,f_k)$ une suite de Kronecker attachée à $\fa$, avec
$f_i\in\A[X^{(i)}]$ ; on pose $\B=\A[X^{(1)},\dots,X^{(k)}]$, $E_\B=E[X]$
et $E_i=E_\B/(f_1,\dots,f_i)E_\B$.

\smallskip\noindent\emph{$\Gr\Rightarrow\Pr$ (9.6 et 9.4.2 seulement).}
Supposons $(f_1,\dots,f_k)$ régulière pour $E_\B$. Comme
$f_k\in\fa\B$ est $E_{k-1}$-régulier,
$H^0(\ua,E_{k-1})=\Ann_{E_{k-1}}(\fa)\subseteq\Ann_{E_{k-1}}(f_k)=0$,
donc $\Pr_\B(\ua,E_{k-1})\ge1$. En remontant avec 9.6 (appliqué à
$f_{k-1},\dots,f_1$), on obtient $\Pr_\B(\ua,E_{k-2})\ge2$, \dots,
$\Pr_\B(\ua,E_\B)\ge k$. Comme $\B$ est libre, donc fidèlement plate,
9.4.2 donne $\Pr_\A(\ua,E)\ge k$.

Cet argument s'applique à n'importe quelle suite $E_C$-régulière contenue
dans $\fa C$, pour $C$ fidèlement plate sur $\A$.

\smallskip\noindent\emph{$\Pr\Rightarrow\Gr$ (9.4.1, McCoy, 9.6).}
Par 9.4.1, $\Pr_\B(\ua,E_\B)\ge k$. Montrons par récurrence sur $i<k$ que
$\Pr_\B(\ua,E_i)\ge k-i$ et que $f_{i+1}$ est $E_i$-régulier. Soit
$\B_i=\A[X^{(j)}:j\ne i+1]$ et
$M_i=E[X^{(j)}:j\neq i+1]/(f_1,\dots,f_i)$ ; puisque $f_1,\dots,f_i$ ne
contiennent pas $X^{(i+1)}$, on a $E_i=M_i[X^{(i+1)}]$ et
$M_i\subseteq E_i$. De $\Pr_\B(\ua,E_i)\ge k-i\ge1$ on tire
$\Ann_{E_i}(\fa)=0$, donc $\Ann_{M_i}(\fa\B_i)=0$. Le contenu de
$f_{i+1}$ vu comme polynôme en $X^{(i+1)}$ sur $\B_i$ est $\fa\B_i$, donc
par le lemme~\ref{mccoy}, $f_{i+1}$ est $E_i$-régulier. Si $k-i-1\ge1$,
9.6 donne $\Pr_\B(\ua,E_{i+1})\ge k-i-1$ et l'on continue.
\end{proof}

\begin{corollary}
Les résultats suivants de la section 1 se déduisent de 9.7 et des
propriétés de l'homologie de Koszul :
\begin{itemize}
\item 1.4 : par la remarque suivant $\Gr\Rightarrow\Pr$ et 9.7 ;
\item 1.13 : de 9.4 ;
\item 1.14 : de 9.6 ;
\item 1.20 : c'est l'exercice 1 (suite exacte longue) ;
\item 1.6 (et la question « Preuve directe en exercice ? » après 9.7) :
  si $\fa\subseteq\fb$ et $\Gr(\fa,E)\ge k$, une suite de Kronecker pour
  $\fa$ est une suite régulière contenue dans $\fb\B$, donc
  $\Pr(\fb,E)\ge k$ par l'argument $\Gr\Rightarrow\Pr$.
\end{itemize}
\end{corollary}

\begin{remark}
On peut aussi voir directement que $\Pr$ ne dépend que de l'idéal : si
$b\in\fa$, la multiplication par $b$ est homotope à $0$ sur
$K^\bullet(\ua,E)$, donc le cône donne
$H^j(\ua,b;E)\simeq H^j(\ua;E)\oplus H^{j-1}(\ua;E)$.
\end{remark}

\noindent On suggère donc de présenter 9.7 non comme une comparaison avec
la théorie de RLF, mais comme un fondement de celle-ci.

\section{Le corolaire 9.8}

Avec $\Pr$ l'énoncé devient transparent, à condition d'ajouter l'hypothèse
« \emph{$E$ de présentation finie} ».

\begin{proposition}
Soit $\A$ cohérent fortement discret et $E$ de présentation finie. Pour
toute suite $\ua$ de longueur $n$, on peut calculer $\Gr_\A(\ua,E)$,
élément de $\iv{0..n}\cup\{\infty\}$.
\end{proposition}
\begin{proof}
Les modules $K^j(\ua,E)=E^{\binom nj}$ sont de présentation finie sur
$\A$ cohérent, donc les $\Ker d_j$ sont de type fini, avec des générateurs
calculables. La forte discrétion permet de tester
$\Ker d_j\subseteq\Ima d_{j-1}$, i.e.\ $H^j=0$. La profondeur est le premier
$j\le n$ tel que $H^j\ne0$, ou $\infty$ si tous les $H^j$ sont nuls (on
conclut par 9.7).
\end{proof}

\noindent Les hypothèses « si la profondeur est finie » et « non trivial »
deviennent inutiles.

\emph{Pourquoi c'était obscur dans RLF.} $\Gr$ est défini dans $\A[X]$, et
la cohérence ne passe pas à $\A[X]$ (contre-exemple de Soublin). Le premier
test (« $f_1$ est-il régulier ? ») se ramène à $\A$ par McCoy, mais les
suivants portent sur $E[X]/f_1E[X]$. Avec $\Pr$, tout reste dans $\A$.

\section{Démonstration du théorème 9.10}

On utilise l'homologie descendante, $\ua'=(a_2,\dots,a_n)$, et la suite
exacte longue (+) :
\[
\cdots\to H_i(\ua';E)\xrightarrow{\times a_1}H_i(\ua';E)\to H_i(\ua;E)\to
H_{i-1}(\ua';E)\xrightarrow{\times a_1}\cdots
\]
ainsi que la version descendante du lemme 9.2, $\fa\, H_i(\ua,E)=0$ (même
homotopie : la multiplication par $a_i$ est homotope à $0$ via la
contraction par $e_i^\star$).

\begin{theorem}[9.10.1]
Soient $\A$ cohérent, $E$ de présentation finie et $k\ge0$. Si
$H_k(\ua,E)=0$, alors $H_i(\ua,E)=0$ pour tout $i\ge k$.
\end{theorem}

\begin{proof}
Récurrence sur $n$, l'énoncé portant sur tous les anneaux cohérents. Pour
$n=0$, le complexe se réduit à $E$ en degré $0$ et l'énoncé est trivial.

Soit $n\ge1$ et $H_k(\ua;E)=0$. Par (+), la multiplication par $a_1$ est
surjective sur $H_k(\ua';E)$, d'où $H_k(\ua';E)=a_1H_k(\ua';E)$.

\emph{Cohérence.} $H_k(\ua';E)$ est de type fini, car $\Ker\partial_k$
est un noyau d'une application entre modules de présentation finie sur un
anneau cohérent. (Pour $k=0$ il suffit que $E$ soit de type fini.)

Le lemme du déterminant fournit $y\in\A$ tel que $s=1+a_1y$ vérifie
$s\,H_k(\ua';E)=0$. Les éléments $s$ et $a_1$ sont comaximaux.

\emph{Dans $\A[1/a_1]$.} On a $\fa=\langle1\rangle$, donc tous les
$H_i(\ua,E)[1/a_1]$ sont nuls (lemme 9.2).

\emph{Dans $\A[1/s]$.} L'anneau $\A[1/s]$ est cohérent et $E_s$ de
présentation finie. Par 9.4, $H_k(\ua',E_s)=H_k(\ua',E)_s=0$. L'hypothèse
de récurrence appliquée à la suite $\ua'$ de longueur $n-1$ donne
$H_i(\ua',E_s)=0$ pour tout $i\ge k$. Pour $i\ge k+1$, la suite (+)
contient $H_i(\ua',E_s)\to H_i(\ua,E_s)\to H_{i-1}(\ua',E_s)$, dont les
deux bouts sont nuls ; donc $H_i(\ua,E)_s=0$.

\emph{Recollement.} Pour $i\ge k+1$, $H_i(\ua,E)$ est nul après
localisation en $s$ et en $a_1$, comaximaux ; par le principe local-global
de base, $H_i(\ua,E)=0$. Le cas $i=k$ est l'hypothèse.
\end{proof}

\noindent Cette démonstration n'utilise ni l'astuce $(a,b,ab)$ (9.9) ni la
théorie de la profondeur ; la cohérence ne sert qu'à garantir que
$H_k(\ua',E)$ est de type fini.

\begin{corollary}[9.10.2]
Si $\ua$ est $1$-sécante pour $E$, elle est complètement sécante pour $E$.
\end{corollary}
\begin{proof}
Par la définition 5.2, $\ua$ est $1$-sécante pour $E$ si et seulement si
$H_1(\ua,E)=0$, puisque $\Ima\partial_2$ est le module des syzygies
triviales. Le théorème donne $H_i=0$ pour tout $i\ge1$. L'autodualité du
complexe de Koszul, $H^j\simeq H_{n-j}$, donne $H^j=0$ pour $j\le n-1$,
donc $\Pr\ge n$, donc $\Gr\ge n$ par 9.7.
\end{proof}
\noindent C'est la réciproque partielle de 5.5 annoncée après celle-ci.
(Dans 1.7.1, « $\Gr\ge n$ » devrait se lire « $\Gr\ge k$ », la suite y
étant notée $a_1,\dots,a_k$.)

\begin{corollary}[9.10.3]
Si $\fa E=E$, alors tous les $H_i(\ua,E)$ sont nuls et
$\Pr(\ua,E)=\Gr(\ua,E)=\infty$.
\end{corollary}
\begin{proof}
$H_0=E/\fa E=0$ ; on conclut par le théorème, ou directement par 9.11 avec
$(1+x)E=0$, $x\in\fa$ (section 1.2).
\end{proof}

\begin{remark}[le « sibrouillon » gradué]
Il faut lire « $H_i=0$ pour $i\ge k$ » en indices descendants (ou
« $i\le k$ » pour $H^\bullet$). La même récurrence marche sans cohérence ni
localisation : on remplace le lemme du déterminant par le Nakayama gradué.
Si $a_1$ est homogène de degré $>0$ et les degrés de $E$ bornés
inférieurement, alors $H=a_1H$ entraîne $H=0$.
\end{remark}

\end{document}
